\documentclass[10pt]{article}
\usepackage[margin=20mm]{geometry}
\usepackage{amsmath,amssymb,fontspec,unicode-math,xurl,hyperref}
\IfFontExistsTF{DejaVu Serif}{\setmainfont{DejaVu Serif}}{\setmainfont{Latin Modern Roman}}
\setmathfont{latinmodern-math.otf}
\setlength{\parindent}{0pt}\setlength{\parskip}{6pt}\sloppy
\title{0062 - Scientific Verification Report}
\author{AI-assisted scientific verification; human decision pending}
\date{24 October 2026}
\begin{document}
\maketitle
No Issue Found within the inspected scope, with scope clarifications. No demonstrated substantive defect supports Major or Minor. Scientific decision remains pending.\par
\section*{1. Identity, version and coverage}
Reviewed paper: Position: In-House Evaluation Is Not Enough. Towards Robust Third-Party Evaluation and Flaw Disclosure for General-Purpose AI, by Shayne Longpre and 33 coauthors. The reviewed version is the 31-page publisher-linked PMLR proceedings PDF. Page references use its printed pages 1-31.\par
SHA-256: 921e087197d00bd2d15bd5a9ee6dcf6fbaf0cbb3cc9a9d849f5fe5ce3bc765b0; size 5959117 bytes. The digest matches the source-acquisition record. The first-page title, author list and proceedings identity were compared with official publisher metadata. The original bytes were preserved.\par
Coverage includes Sections 1-7 (pp.1-9), definitions in Appendix A (pp.18-20), report schemas and examples in Appendix B (pp.20-25), policy and safe-harbor qualifications in Appendix C (pp.25-28), and the reporting/taxonomy inventory in Appendix D (pp.29-31). Selected primary support was checked; this is not a verification of every historical law, product policy or reference. The independent reviewer inspected the p.21 case illustrations; the root reviewer visually inspected p.26 Figure A7.\par
\section*{2. Core contribution and checked dependencies}
The paper proposes standardized flaw reports and conduct rules, provider legal/technical protections, and a coordinating center. It does not present a new statistical estimator, a formal security proof or a deployed implementation with measured performance. The review tested the proposal’s stated dependencies, the factual examples used to motivate it, and whether qualifications in the body resolve apparent overbroad readings.\par
Claim / location | Check and outcome\par
Third-party scrutiny and cross-provider coordination; Sections 2-3, pp.3-4; Section 4.3, pp.7-8 | The cited cross-model transfer premise has direct support in Zou et al., arXiv:2307.15043v2. This supports the existence of a coordination need; it does not prove that every flaw transfers or that a particular hub will be effective.\par
Traceable, reproducible reports; Section 4.1, p.5; Table A1, p.23 | The schema records system versions, context, timestamps, session identifiers and the claimed violation; it also includes replication material and statistical arguments for appropriate report types. The text recommends frequency evidence for probabilistic outputs. No one-shot failure was represented as a measured population failure rate.\par
Legal and technical safe harbors; Section 4.2, pp.6-7; C.3, pp.26-28 | The provider template limits commitments to its own conduct, requires compliance with applicable law, and separately addresses third-party litigation. Moderation-exempt access requires vetting, and the text discusses erroneous suspension and appeals. Figure A7 is a schematic of this proposed access regime, not an unconditional immunity theorem.\par
495-person OLMo event; B.1, p.20 | The original McGregor et al. case-study record corroborates 495 participants and documentation-focused adjudication. This is a cited historical pilot, not a new experiment performed by the present paper.\par
Example flaw cards; B.1, pp.20-21 | The cards are explicitly illustrative reconstructions of earlier training-data extraction and gender-bias reports. The BERT example corresponds to AVID-2022-R0001, version 0.1, reported 2022-11-09. Their existence supports the report-schema illustration, not the efficacy of the proposed national or cross-provider infrastructure.\par
Disclosure timing; Figure 1, p.2; Section 4.1 footnote 5, p.6; Section 4.3, p.8 | The diagram’s approximate 30-day label is qualified by harm-dependent periods, exceptions and possible extensions in the text. Treating it as a universal compulsory deadline would misrepresent the proposal.\par
Resource, governance and remediation questions; Sections 5-6, pp.8-9 | The text discusses resource tradeoffs and explicitly leaves governance, responsibility disputes and effective remediation for future work. These are open implementation/evaluation tasks, not concealed experimental results.\par
\section*{3. Scope clarifications and material uncertainties}
The 2022 DOJ policy is an instruction about federal charging decisions; its final paragraph disclaims creating an enforceable right. It is not blanket immunity. Read together with the paper’s provider-specific template and call for further legal changes (C.1, p.25), the paper’s use of this precedent does not establish a contradiction in the core proposal. Figure A7’s shorthand should be interpreted within those express limits.\par
The paper’s extension from software disclosure practice to broad sociotechnical AI flaws is an argued policy/design recommendation. It is not established as an experimentally measured causal improvement in this article. The claimed low cost and future benefits have not been quantified here. An empirical rollout could assess these claims; an absent rollout is not, by itself, an error in a normative recommendation.\par
The proposed tag-routing center is an architectural sketch. Subscriber authentication, confidential-report access control, versioned policy handling, governance and operational effectiveness would require specification and testing before deployment. The inspected text makes no complete implementation-security guarantee that can be refuted by imagining an unspecified insecure implementation.\par
No author-code execution was performed. No new evaluation-code repository or independent experimental dataset was identified in the inspected paper/publisher materials. Linked standards, report databases and prior studies were treated as supporting sources, not as implementations of this paper.\par
Case-source verification did not rerun the original attacks, recount raw event participation logs, or independently recompute the approximately 3\% frequency and longest-sequence detail printed in the illustrative extraction card. Its example timestamp was not treated as an authenticated historical disclosure date.\par
\section*{4. Overall outcome and remaining work}
Item | Judgment | Repair / adoption | Remaining work\par
Core disclosure framework | No Issue Found in inspected reasoning | No correction established; adoption not requested | Implementation and effectiveness evaluation.\par
Historical case examples | No Issue Found in checked attribution | Not applicable | Underlying attack and bias experiments not rerun.\par
Safe-harbor scope | Scope clarification | No verified repair needed | Jurisdiction-specific legal implementation is outside this scientific check.\par
Scientific decision | Pending | Human review not asserted | publication status and independent human assessment.\par
The arXiv record for 2503.16861 shows v2 dated25 March 2025. That preprint was used for version discovery, not substituted for the publisher PDF. No separate correction was identified in the focused search. A comprehensive comparison of all preprint/OpenReview revisions was not performed.\par
No Problem section is manufactured because this review did not demonstrate a substantive defect. This scoped result does not certify every policy inventory entry, reference, future implementation or the empirical success of the recommendations.\par
\section*{5. Consulted references}
Longpre et al., official ICML 2025 proceedings record: https://proceedings.mlr.press/v267/longpre25a.html\par
Longpre et al., preprint version history (v2, 25 March 2025): https://arxiv.org/abs/2503.16861\par
Zou et al., Universal and Transferable Adversarial Attacks on Aligned Language Models, v2, 20 December 2023 (abstract-level premise check): https://arxiv.org/abs/2307.15043v2\par
McGregor et al., To Err is AI, arXiv:2410.12104v1, 15 October 2024 (case-study record): https://arxiv.org/abs/2410.12104\par
AVID gender-bias report, AVID-2022-R0001, version 0.1: https://avidml.org/database/avid-2022-r0001/\par
US DOJ, CFAA charging policy, 19 May 2022; Section B.3 and final limitation: https://www.justice.gov/d9/press-releases/attachments/2022/05/19/cfaa\_policy\_may\_19\_0.pdf\par
Nasr et al., Scalable Extraction of Training Data from (Production) Language Models, arXiv:2311.17035v1, 28 November 2023 (selected case-source inspection): https://arxiv.org/abs/2311.17035v1\par
\section*{Supporting checks}
The exact disclose.io template URL printed in the paper returned HTTP404 in the current retrieval, so an external template-to-template comparison is not claimed. The review instead inspected the complete provider template included on p.7. This broken retrieval does not establish a scientific defect or absence of a historical template.\par
No author-code execution occurred.\par
No human scientific verification or author adoption is asserted.\par
Human scientific assessment remains pending.\par
\end{document}
